\documentclass[12pt,a4paper]{article}
\usepackage[utf8]{inputenc}
\usepackage[T1]{fontenc}
\usepackage[brazil]{babel}
\usepackage{amsmath,amssymb}
\usepackage{geometry}
\geometry{margin=2.5cm}
\title{Material de Estudo: Dimens\~ao, Escala e Fractais}
\author{Princ\'ipio de Modelagem Matem\'atica}
\date{Unidade 05}
\begin{document}
\maketitle

\section*{Objetivo}
Compreender o conceito de dimens\~ao fractal, calcular dimens\~oes de fractais cl\'assicos e reconhecer padr\~oes fractais em fen\^omenos naturais.

\section*{Conceitos Fundamentais}

\subsection*{Dimens\~ao Topol\'ogica vs.\ Fractal}
A dimens\~ao topol\'ogica \'e um inteiro (ponto=0, linha=1, plano=2, volume=3).
A \textbf{dimens\~ao fractal} $d_f$ pode ser n\~ao-inteira e mede como a complexidade de um objeto cresce com a resolu\c{c}\~ao.

\subsection*{F\'ormula de Hausdorff-Besicovitch}
Para um fractal auto-similar com $N$ c\'opias menores, cada uma reduzida por fator $1/r$:
\[
d_f = \frac{\log N}{\log r}.
\]

\subsection*{Fractais Cl\'assicos}
\begin{itemize}
  \item \textbf{Floco de Koch:} $N=4$, $r=3$ $\Rightarrow$ $d_f = \log 4 / \log 3 \approx 1{,}26$.
  \item \textbf{Tri\^angulo de Sierpinski:} $N=3$, $r=2$ $\Rightarrow$ $d_f = \log 3 / \log 2 \approx 1{,}585$.
  \item \textbf{Esponja de Menger:} $N=20$, $r=3$ $\Rightarrow$ $d_f = \log 20 / \log 3 \approx 2{,}727$.
  \item \textbf{Conjunto de Cantor:} $N=2$, $r=3$ $\Rightarrow$ $d_f = \log 2 / \log 3 \approx 0{,}631$.
\end{itemize}

\subsection*{M\'etodo de Contagem de Caixas}
Para objetos experimentais: cubra com caixas de lado $\varepsilon$ e conte $N(\varepsilon)$. O gr\'afico $\log N$ vs.\ $\log(1/\varepsilon)$ tem inclina\c{c}\~ao $d_f$.

\subsection*{Fractais na Natureza}
Costas litor\^aneas ($d_f \approx 1{,}2$--$1{,}3$), bronqu\'iolos pulmonares, redes de rios, rel\^ampagos, flocos de neve, Conjunto de Mandelbrot.

\section*{Exemplos Resolvidos}

\subsection*{Exemplo 1 --- Curva de Koch: per\'imetro infinito}
A cada itera\c{c}\~ao, o comprimento multiplica por $4/3$. Ap\'os $n$ itera\c{c}\~oes:
\[
L_n = L_0 \left(\frac{4}{3}\right)^n \xrightarrow{n\to\infty} \infty.
\]
O per\'imetro \'e infinito, mas o fractal cabe dentro de um tri\^angulo finito! (\'Area converge.)

\subsection*{Exemplo 2 --- Dimens\~ao da costa brit\^anica (Richardson)}
Richardson (1961) mediu a costa brit\^anica com r\'eguas de comprimento $\varepsilon$. Resultado:
\[
L(\varepsilon) \approx k\,\varepsilon^{1-d_f}, \quad d_f \approx 1{,}25.
\]
Quanto menor a r\'egua, mais detalhes da costa s\~ao capturados e maior o comprimento medido.

\subsection*{Exemplo 3 --- C\'alculo da dimens\~ao de Sierpinski}
\begin{itemize}
  \item Itera\c{c}\~ao 0: 1 tri\^angulo de lado 1.
  \item Itera\c{c}\~ao 1: 3 tri\^angulos de lado 1/2.
  \item Itera\c{c}\~ao 2: 9 tri\^angulos de lado 1/4.
  \item Itera\c{c}\~ao $n$: $3^n$ tri\^angulos de lado $(1/2)^n$.
\end{itemize}
\[
d_f = \frac{\log 3^n}{\log 2^n} = \frac{\log 3}{\log 2} \approx 1{,}585.
\]

\section*{Exerc\'icios}

\begin{enumerate}
  \item \textbf{(Dimens\~ao)} Calcule a dimens\~ao fractal de cada objeto:
    \begin{enumerate}
      \item[$a.$] Tapete de Sierpinski: $N=8$ c\'opias, fator $r=3$.
      \item[$b.$] Fractal com $N=5$ c\'opias, fator $r=5$.
      \item[$c.$] Conjunto de Cantor: $N=2$, $r=3$. Interprete o resultado ($d_f < 1$).
    \end{enumerate}
  
  \item \textbf{(Koch)} Na curva de Koch, quantas itera\c{c}\~oes s\~ao necess\'arias para que o comprimento seja 10 vezes o original? (Equa\c{c}\~ao: $L_n = L_0(4/3)^n$.)
  
  \item \textbf{(Contagem de caixas)} Para um objeto fractal, os dados medidos s\~ao:
    \begin{center}
    \begin{tabular}{|c|c|c|c|}
    \hline
    $1/\varepsilon$ & 10 & 100 & 1000 \\ \hline
    $N(\varepsilon)$ & 45 & 1350 & 42\,500 \\ \hline
    \end{tabular}
    \end{center}
    Estime $d_f$ usando os dados.
  
  \item \textbf{(\'Area da costa)} Se $d_f = 1{,}25$ para a costa brit\^anica e $k = 2000\,\text{km}^{d_f}$, qual o comprimento medido com uma r\'egua de 1 km? E com uma de 100 m? Comente.
  
  \item \textbf{(Pulmao)} Os bronqu\'iolos pulmonares s\~ao fractais com $d_f \approx 2{,}8$. Por que uma dimens\~ao fractal maior que 2 \'e vantajosa para a troca de gases?
  
  \item \textbf{(Mandelbrot)} O conjunto de Mandelbrot \'e definido por $z_{n+1} = z_n^2 + c$ para $z_0=0$, $c \in \mathbb{C}$. Um ponto $c$ est\'a no conjunto se $|z_n|$ n\~ao diverge. Verifique para: (a) $c=0$; (b) $c=1$; (c) $c=-1$.
  
  \item \textbf{(Auto-similaridade)} Explique a diferen\c{c}a entre auto-similaridade \emph{exata} (fractais matem\'aticos) e \emph{estat\'istica} (fractais naturais). D\^e um exemplo de cada.
\end{enumerate}

\section*{Resumo R\'apido}
\begin{itemize}
  \item Dimens\~ao fractal: $d_f = \log N / \log r$ (auto-similares) ou inclina\c{c}\~ao do gr\'afico log-log.
  \item Koch ($d_f\approx1{,}26$), Sierpinski ($d_f\approx1{,}585$), Cantor ($d_f\approx0{,}631$).
  \item Fractais modelam costas, pulm\~oes, redes de rios, rel\^ampagos.
\end{itemize}

\section*{Refer\^encias}
\begin{itemize}
  \item Mandelbrot, B.\ B.\ \textit{The Fractal Geometry of Nature}. Freeman, 1982.
  \item Peitgen, H.-O.\ et al.\ \textit{Chaos and Fractals}. Springer, 2004.
  \item Notas de aula da disciplina.
\end{itemize}

\end{document}
